\section{Detailed P3 typed-role extraction and port geometry}
\label{app:p3technical}

This appendix preserves the complete P3 extraction argument while keeping the
main proof spine visible.  It contains Propositions~19--31: the squared-Hodge
realization, the complete first-jet quotient, source-rank and typed-extraction
tests, represented-role and common-scale candidates, the O/Q and M/G
stabilizer analysis, and the cotangent-port construction.  None of these
technical reductions proves that the unrestricted physical base--seed parent
occupies the required complete typed incidence.

\begin{proposition}[Full-role squared-Hodge selector and P3 realization]
\label{prop:p3clifford}
Let $E_{P3}$ be the seven-role physical quotient with positive metric $K_*$.
If one source-owned linear self-adjoint operator map
$\mathsf C:E_{P3}\to\mathfrak A_{\mathrm{src}}$ satisfies
\begin{equation}
\mathsf C(v)^2=K_*(v,v)I
\qquad\text{for every }v\in E_{P3},
\label{eq:fullrolesquare}
\end{equation}
then
\begin{equation}
\mathsf C(u)\mathsf C(v)+\mathsf C(v)\mathsf C(u)
=2K_*(u,v)I.
\label{eq:p3clifford}
\end{equation}
Thus an orthonormal seven-role frame generates
$\operatorname{Cl}_7(\mathbb C)\simeq
M_8(\mathbb C)\oplus M_8(\mathbb C)$.  In the canonical eight-qubit
Jordan--Wigner embedding into the common
$\operatorname{Cl}_{16}(\mathbb C)\simeq M_{256}(\mathbb C)$ envelope, the
two simple sectors occur with multiplicities $(16,16)$ and normalized-trace
weights $(1/2,1/2)$.
\end{proposition}

\begin{proof}
Apply Eq.~\eqref{eq:fullrolesquare} to $u+v$ and subtract its $u$ and $v$
instances to obtain Eq.~\eqref{eq:p3clifford}.  The complex Clifford
classification gives the two simple rank-eight sectors.  The canonical first
seven Jordan--Wigner generators have 128 linearly independent Clifford words
and central volume eigenvalues of opposite sign with equal eight-dimensional
multiplicity in the minimal faithful direct sum.  Sixteen copies embed this
faithful representation into dimension 256, giving the stated multiplicities
and trace weights.  These finite checks are regenerated by
\path{scripts/audit_theorem_certificates.py}.
\end{proof}

The operator-square premise is stronger than a normalized trace-Gram or
scalar squared-action identity.  Seven commuting diagonal matrices can have
the same identity trace Gram as seven Clifford generators while their generic
linear combination does not square to a scalar.  Hence the P3 packet is an
explicit coefficient-free completion and Eq.~\eqref{eq:fullrolesquare} is an
exact selector, but the unrestricted base--seed reduct has not selected the
required typed seven-role operator incidence as part of one occupied packet.
The existing rank-two EACER and rank-nine
seed/context squared-Hodge carriers have the wrong domains.  A post-body
Clifford handoff is therefore an admissible minimal working completion, not an
unrestricted base--seed derivation or a modification of the already stabilized body.

There is a further type distinction.  The seven entries in the P3 support
ledger are unital role algebras, not seven canonically one-dimensional source
slots.  Their units are identified in the unital primitive colimit and hence
span only one common line.  Selecting one non-scalar element inside each role
algebra would add noncanonical data and need not be invariant under internal
role automorphisms.  Therefore the support ledger and the current colimit do
not by themselves construct the seven-dimensional domain used in
Proposition~\ref{prop:p3clifford}.

The sufficient source object can be stated without choosing such generators.
Let $(\mathcal P_{P3}^1,G)$ be the complete de-duplicated real physical
first-jet module after the true-null quotient, and suppose the source owns a
surjection
\begin{equation}
q_{P3}:\mathcal P_{P3}^1\longrightarrow E_{P3}^{\rm lbl}\simeq\mathbb R^7
\label{eq:p3rolequotient}
\end{equation}
whose seven typed rows are nonzero, whose kernel is annihilated by every
claimed P3 first-jet terminal, and through which every relevant degree-one
P3 response factors.  If the common action unit normalizes it as
\begin{equation}
q_{P3}G^{-1}q_{P3}^{*}=I_7,
\label{eq:p3coisometry}
\end{equation}
then its metric adjoint
$j_{P3}=G^{-1}q_{P3}^{*}$ is a canonical isometric section:
\begin{equation}
q_{P3}j_{P3}=I_7,
\qquad
j_{P3}^{*}G j_{P3}=I_7.
\label{eq:p3isometricsection}
\end{equation}
The exterior construction on $E_{P3}^{\rm lbl}$ then supplies the normalized
Clifford law of Proposition~\ref{prop:p3clifford}.  This is a conditional
role-quotient theorem, not a derivation of Eq.~\eqref{eq:p3rolequotient} from
the physical base--seed system.  Existing seed/body decoders, primitive-edge
label algebras and downstream separating-record witnesses have different
typed domains and cannot be reused as $q_{P3}$ without an additional source
theorem.

The normalization in Eq.~\eqref{eq:p3coisometry} is not an independent
physical parent-law coefficient if a complete typed role differential is already
source-owned.  For a surjection
$R_{P3}:\mathcal P_{P3}^1\to E_{P3}^{\rm lbl}$, define
\begin{equation}
H_R=R_{P3}G^{-1}R_{P3}^{*}>0,
\qquad
q_{P3}=H_R^{-1/2}R_{P3}.
\label{eq:p3polarquotient}
\end{equation}
The unique positive square root gives
$q_{P3}G^{-1}q_{P3}^{*}=I_7$ and
$\ker q_{P3}=\ker R_{P3}$.

The kernel and no-bypass clauses can be reduced further.  Let
$\{L_\alpha\}_{\alpha\in\mathfrak F_{P3}^{1}}$ be the complete declared
family of degree-one P3 operator and terminal jets on
$\mathcal P_{P3}^{1}$, and set
\begin{equation}
L_\oplus=\bigoplus_\alpha L_\alpha,
\qquad
N_{P3}^{1}=\ker L_\oplus,
\qquad
E_{P3}^{\rm phys}=\mathcal P_{P3}^{1}/N_{P3}^{1}.
\label{eq:p3nullquotient}
\end{equation}

\begin{proposition}[Canonical complete-first-jet quotient]
\label{prop:p3nullquotient}
With the minimum-lift metric induced by $G$, the canonical quotient map
$q_{\rm null}:\mathcal P_{P3}^{1}\to E_{P3}^{\rm phys}$ is surjective,
has kernel $N_{P3}^{1}$, and every $L_\alpha$ factors uniquely through it.
Its horizontal metric adjoint is the minimum-norm isometric section.  Any
other coisometric quotient with the same kernel differs only by an isometry
of $E_{P3}^{\rm phys}$, restricted to a typing-preserving isometry when the
role labels are retained.  Moreover,
\begin{equation}
\dim E_{P3}^{\rm phys}=\operatorname{rank}L_\oplus.
\label{eq:p3nullquotientrank}
\end{equation}
Consequently a complete rank-seven first-jet family supplies the carrier used
in Proposition~\ref{prop:p3clifford} without an additional role-decoder,
normalization or degree-one-bypass choice.
\end{proposition}

\begin{proof}
The first isomorphism theorem gives the kernel, surjectivity and
Eq.~\eqref{eq:p3nullquotientrank}.  Since
$N_{P3}^{1}\subseteq\ker L_\alpha$ for every $\alpha$, each map is constant
on quotient fibres and therefore factors uniquely.  Positivity of $G$ gives
the orthogonal splitting
$\mathcal P_{P3}^{1}=N_{P3}^{1}\oplus(N_{P3}^{1})^{\perp_G}$; restriction of
the quotient map to the horizontal summand defines the minimum-lift metric and
the isometric section.  Two coisometries with the same kernel restrict to
isometries from this same horizontal summand, so their relative map is an
isometry.  QED.
\end{proof}

If a proposed additional degree-one row does not annihilate $N_{P3}^{1}$,
then it does not constitute an independent bypass parameter of the same
complete quotient: it strictly refines the joint null and proves that the
declared family was incomplete.  Thus the remaining upstream premise is only
that the physical base--seed law supplies the complete family in
Eq.~\eqref{eq:p3nullquotient} with rank seven.  The current seed-presentation
candidate has the abstract ingredients from which such a family could be
formed, but the present seed grammar proves neither its completeness nor its
rank in the physically realized packet.

This upstream qualification does not reopen the three downstream levels
already assembled in the enriched MVP: common-source terminal factorization,
the occupied P3 algebra--record--GNS packet, and the conditional terminal
recoveries remain intact under their declared assumptions.  The unresolved
statement is only that the unrestricted physical base--seed parent must select
that joint packet.

The rank premise has an exact finite source certificate.  Let
$K_{\rm src}>0$ be the frozen source Hessian on the true-null-quotiented
first-jet module and let $L_\oplus$ have seven typed rows.  Define
\begin{equation}
G_{P3}^{(1)}=L_\oplus K_{\rm src}^{-1}L_\oplus^*.
\label{eq:p3rankgram}
\end{equation}

\begin{proposition}[Rank-seven source certificate and current-parent no-go]
\label{prop:p3rankcertificate}
The following statements are equivalent:
\begin{equation}
\operatorname{rank}L_\oplus=7
\quad\Longleftrightarrow\quad
G_{P3}^{(1)}>0
\quad\Longleftrightarrow\quad
\lambda_{\min}(G_{P3}^{(1)})>0
\quad\Longleftrightarrow\quad
\det G_{P3}^{(1)}>0.
\label{eq:p3rankequivalence}
\end{equation}
Strict source stability, nonzero activation of all seven typed rows, equal
primitive row costs and a common action unit do not imply these conditions.
Consequently the conditional three-level closure cannot be inverted to prove
rank seven.
\end{proposition}

\begin{proof}
Equation~\eqref{eq:p3rankgram} is the Gram matrix of
$L_\oplus K_{\rm src}^{-1/2}$, and $K_{\rm src}^{-1/2}$ is invertible.
Therefore its rank equals $\operatorname{rank}L_\oplus$; the remaining
equivalences are the finite-dimensional full-rank criteria for a positive
semidefinite $7\times7$ matrix.  For non-entailment, take the same source
space, $K_{\rm src}=I_7$, load and stationary point.  Let $L_7=I_7$, whereas
$L_6$ has rows $e_1,\ldots,e_6$ and
$(e_1+\cdots+e_6)/\sqrt6$.  Both maps have seven nonzero unit rows and total
row-Gram trace seven, but $\operatorname{rank}L_7=7$ and
$\operatorname{rank}L_6=6$.  Hence all the listed upstream properties hold
in both controls while Eq.~\eqref{eq:p3rankequivalence} holds only in the
first.  QED.
\end{proof}

The PS1 triangular dependency action already supplies the untyped rank on its
own codomain.  Let $\pi_{\rm dep}$ delete the root residual and define
\begin{equation}
D_{\rm dep}=\pi_{\rm dep}D_\Sigma:
\bigoplus_{i=0}^{7}E_i\longrightarrow
Z_{\rm dep}:=\bigoplus_{i=1}^{7}E_i.
\label{eq:p3dependencyoperator}
\end{equation}

\begin{proposition}[Onto dependency response and typed-extraction reduction]
\label{prop:p3typedextraction}
The map $D_{\rm dep}$ is surjective and
\begin{equation}
D_{\rm dep}H_{\rm src}^{-1}D_{\rm dep}^*
=\mathcal A_*^{-1}\pi_{\rm dep}K_\Sigma^{-1}\pi_{\rm dep}^*>0.
\label{eq:p3dependencygram}
\end{equation}
If a source-owned typed extraction
$T_{P3}:Z_{\rm dep}\to E_{P3}^{\rm lbl}$ realizes the complete P3 first-jet
family by $L_\oplus=T_{P3}D_{\rm dep}$, then
\begin{equation}
\operatorname{rank}L_\oplus=\operatorname{rank}T_{P3}.
\label{eq:p3typedrank}
\end{equation}
Thus rank seven follows exactly when this complete typed extraction is onto.
The present PS1 action does not select that extraction.
\end{proposition}

\begin{proof}
The restriction of $D_{\rm dep}$ to
$\bigoplus_{i=1}^{7}E_i$ is block lower triangular with identity diagonal,
so it is invertible.  Inverting
$H_{\rm src}=\mathcal A_*D_\Sigma^*K_\Sigma D_\Sigma$ and substituting
$D_{\rm dep}=\pi_{\rm dep}D_\Sigma$ gives
Eq.~\eqref{eq:p3dependencygram}; it is positive because it is a principal
compression of $K_\Sigma^{-1}>0$.  Surjectivity gives
$\operatorname{im}(T_{P3}D_{\rm dep})=\operatorname{im}T_{P3}$, proving
Eq.~\eqref{eq:p3typedrank}.  The rank-six and rank-seven extraction controls
in Proposition~\ref{prop:p3rankcertificate} share the same dependency action,
so that action cannot select between them.  QED.
\end{proof}

The seven dependency blocks $S1,\ldots,S7$ are not the seven generally
non-scalar P3 ownership roles $\{B,O,R,U,Q,M,G\}$.  Consequently this result
does not authorize a label-based identification.  It shows instead that raw
seven-direction source response is already available: the remaining
unrestricted physical datum is one sourced, complete and onto P3-typed
extraction.  No new carrier, decoder gain or terminal normalization is
required, and downstream terminal maps may not be used to choose it.

The common represented trace does not by itself construct that extraction.
On a simple role block $\mathcal A_a=M_n(\mathbb C)$, every real linear
functional on $\mathcal A_a^{\rm sa}$ invariant under all inner
automorphisms is a multiple of the trace and therefore vanishes on the
centered tangent
$\mathcal A_{a,0}^{\rm sa}=\{x=x^*: \omega_a(x)=0\}$.  Thus neither the
trace nor the role ledger canonically chooses a noncentral role direction.

The occupied seed representation nevertheless supplies a canonical
candidate when its role values are noncentral.  Let
$E_a:X_{\rm src}\to\mathcal A_a^{\rm sa}$ be the source-owned represented
role evaluation and define

\begin{equation}
h_a=E_a(x_*)-\omega_a(E_a(x_*))I,
\qquad
\ell_a(\delta x)=
{\operatorname{Re}\omega_a(h_aD E_a\delta x)
 \over \sqrt{\omega_a(h_a^2)}}.
\label{eq:p3radialcovector}
\end{equation}

This is defined for $h_a\ne0$ and is natural under simultaneous unitary
conjugation of the occupied value and its derivative.  Let
$L_{\rm rad}=\bigoplus_a\ell_a:X_{\rm src}\to\mathbb R^7$.

\begin{proposition}[Occupied-role radial extraction criterion]
\label{prop:p3radialextraction}
There is a unique typed map
$T_{P3}^{\rm rad}:Z_{\rm dep}\to E_{P3}^{\rm lbl}$ satisfying
$L_{\rm rad}=T_{P3}^{\rm rad}D_{\rm dep}$ if and only if
\begin{equation}
\ker D_{\rm dep}\subseteq\ker L_{\rm rad}.
\label{eq:p3radialdescent}
\end{equation}
When it exists, it is onto if and only if
\begin{equation}
\operatorname{rank}L_{\rm rad}=7
\quad\Longleftrightarrow\quad
L_{\rm rad}H_{\rm src}^{-1}L_{\rm rad}^{*}>0.
\label{eq:p3radialgram}
\end{equation}
The present source ledger does not entail either condition.
\end{proposition}

\begin{proof}
Surjectivity of $D_{\rm dep}$ identifies $Z_{\rm dep}$ with
$X_{\rm src}/\ker D_{\rm dep}$.  The universal property of this quotient
gives the unique factor precisely under
Eq.~\eqref{eq:p3radialdescent}.  Its image equals the image of
$L_{\rm rad}$, and the final equivalence is the positive Gram criterion of
Proposition~\ref{prop:p3rankcertificate}.  Non-entailment is sharp: a zero
$h_a$ removes a row; two nonzero equal-cost anchors may induce the same
covector and give rank six; and a covector nonzero on $\ker D_{\rm dep}$
does not descend.  QED.
\end{proof}

This replaces an arbitrary $T_{P3}$ by a derivative of occupied role data
whenever the physical seed representation supplies the seven anchors.  It
does not prove that the physical parent law supplies them, that the deleted root is invisible
to all seven radial covectors, or that their Gram is positive definite.  A
non-radial typed differential remains possible only as separately sourced
physics.

The common-unit role requires a typed correction to the all-radial candidate.
If $E_U(x_*)=cI$, then its centered anchor is zero and
Eq.~\eqref{eq:p3radialcovector} supplies no $U$ row.  If the source instead
owns a differentiable positive scale/clock-unit character
$u_*:X_{\rm src}\to\mathbb R_+$, define
\begin{equation}
\kappa_U=d\log u_*|_{x_*},
\qquad
\ell_U^{\rm sc}={\kappa_U\over
\sqrt{\kappa_UH_{\rm src}^{-1}\kappa_U^*}}.
\label{eq:p3unitscalerow}
\end{equation}
Replacing the zero centered $U$ row by $\ell_U^{\rm sc}$ gives a hybrid
family to which Proposition~\ref{prop:p3radialextraction} applies unchanged.
The common action unit normalizes a supplied nonzero unit incidence; it does
not make $\kappa_U$ nonzero.  Hence the current source packet must still
derive that scale/clock differential or another independently typed $U$
incidence.

\begin{proposition}[Current-source metric non-identifiability]
\label{prop:p3currentevaluability}
The current PS1 source ledger does not determine the hybrid seven-role Gram,
even after restricting to full-rank extractions with equal row cost.
\end{proposition}

\begin{proof}
The ledger gives the conditional common tangent $J_*$ and the aggregate
restrictions $e_{P1}J_*$ and $e_OJ_*$, but not the seven represented
evaluation jets $E_a(x_*),DE_a|_{x_*}$ or a selected $d\log u_*$.  This
already prevents evaluation of Eqs.~\eqref{eq:p3radialcovector} and
\eqref{eq:p3unitscalerow}.  Non-identifiability is sharp.  Whiten the fixed
onto dependency response to
$Q_{\rm dep}H_{\rm src}^{-1}Q_{\rm dep}^*=I_7$.  Set $T_0=I_7$ and, for
$0<|c|<1$, let $T_c$ leave all rows fixed except
$e_2^*T_c=ce_1^*+\sqrt{1-c^2}\,e_2^*$.  Both
$L_0=T_0Q_{\rm dep}$ and $L_c=T_cQ_{\rm dep}$ factor through the same
dependency quotient, have rank seven, and have unit diagonal Gram.  However,
the second Gram has off-diagonal entry $c$ and determinant $1-c^2$, whereas
the first is $I_7$.  Therefore the common reduct, rank and row normalization
do not select the physical P3 metric.  QED.
\end{proof}

The finite missing datum is one occupied represented role-evaluation jet
$\{E_a(x_*),DE_a|_{x_*}\}_{a\ne U}$ together with the source-owned unit row
$d\log u_*|_{x_*}$.  Supplying that packet from the physical base--seed
representation would make the descent and positive-Gram tests decisive in
one calculation.  Choosing it from terminal behavior would instead be a
fitted completion.

The common-source terminal assembly cannot be inverted to remove this
missing datum.  If $J_{\rm term}=\mathcal T A_O$ is the stacked terminal
precursor, a factorization $L_{\rm mix}=R J_{\rm term}$ exists exactly when
$\ker J_{\rm term}\subseteq\ker L_{\rm mix}$.  The forward MOPR theorem gives
only $\operatorname{rank}J_{\rm term}\leq\operatorname{rank}A_O$ and does not
establish this kernel inclusion.  Terminal recovery of the source jet would
therefore require an additional source-owned injectivity theorem and cannot
be inferred from successful terminal limits.

The closest existing source module does not remove the gap by dimension
matching.  The conditional HRC role one-jet is
$V_{\rm HRC}=E_S\oplus E_S$, two copies of the three-dimensional standard
$S_4$ representation.  The six non-unit P3 entries, however, are fixed
semantic ownership labels under the current type-preserving ledger; no
tetrahedral action mixing them has been derived.  Therefore
\begin{equation}
\operatorname{Hom}_{S_4}
\left(E_S\oplus E_S,\mathbf1^{\oplus6}\right)=0.
\label{eq:p3hrctypenogo}
\end{equation}
Indeed, each row of an equivariant map would be an invariant covector on a
standard module, and no such nonzero covector exists.  A matched
$E_S\oplus E_S$ target action admits the identity intertwiner, so the
obstruction is typing rather than dimension.  Deriving such a target action
would itself be the missing covariant role bridge and would have to preserve
the distinct typed relation-support classes.  Consequently the count
$6+1=7$, even after adding the independent $U$ scale row, does not construct
the occupied P3 evaluation jet.

The six non-unit maps can nevertheless be reduced to one existing source
object.  Let $\mathcal A_{\rm curr}$ be the de-duplicated primitive colimit,
$j_a:\mathcal A_a\to\mathcal A_{\rm curr}$ its canonical role maps, and let
$x\mapsto\Pi_x$ be a $C^1$ germ of the single coherent physical unital
$*$-homomorphism.  For source-declared self-adjoint anchors
$z_a\in\mathcal A_a$, define
\begin{equation}
E_a(x)=\Pi_x(j_a(z_a)),
\qquad
DE_a|_{x_*}=D\Pi|_{x_*}(j_a(z_a)).
\label{eq:p3coconejet}
\end{equation}
All six non-unit evaluation jets are then restrictions of one representation
one-jet.  Differentiating the homomorphism and dagger identities gives
\begin{equation}
D\Pi(ab)=D\Pi(a)\Pi(b)+\Pi(a)D\Pi(b),
\qquad D\Pi(a^*)=D\Pi(a)^*,
\qquad D\Pi(1)=0.
\label{eq:p3coconelinearized}
\end{equation}
Thus the missing maps are not six independent completion choices once the
physical typed presentation, anchors and cocone germ are supplied.  The last
identity also proves that the algebraic unit cannot generate the nonzero
absolute $U$ scale row.  That row must come from a distinguished positive
$U$-role scale generator or from the separately sourced character $u_*$.  The
remaining finite source datum is therefore one occupied cocone one-jet, its
typed anchor tuple and the unit-scale differential.  The current parent has
not selected this packet, and cocone coherence by itself does not imply the
rank-seven positive-Gram condition.

The unit-scale differential is not wholly structureless on the conditional
protected relation-unit branch already isolated upstream.  If the physical
source occupies the commuting relation-unit square in a source-frozen chart,
then the positive character may be chosen as
$u_*=U_R=a_*/c_R$, and therefore
\begin{equation}
 \boxed{\kappa_U=d\log U_R=d\log a_*-d\log c_R.}
 \label{eq:p3relationscalepullback}
\end{equation}
In the common CHDF Hodge chart, where $U_R=\sqrt2\,\ell_H$, this becomes
$\kappa_U=d\log\ell_H$.  This is a conditional pullback identity, not an
activation theorem.  Equal logarithmic source responses give
$\kappa_U=0$; moreover, the finite common-unit witness fixes ratios at one
packet but does not itself provide a nonzero source tangent.  Thus the
remaining question is physical occupation of the protected square together
with nonzero activation, rather than the form of an otherwise arbitrary
seventh row.

There is also an exact finite novelty certificate.  Let $L_6$ collect the six
descended non-unit rows and assume
$G_6=L_6H_{\rm src}^{-1}L_6^*>0$.  With
$b_U=L_6H_{\rm src}^{-1}\kappa_U^*$, define
\begin{equation}
 \sigma_U=\kappa_UH_{\rm src}^{-1}\kappa_U^*
 -b_U^*G_6^{-1}b_U.
 \label{eq:p3uscaleSchur}
\end{equation}
The block-Gram Schur-complement theorem gives
\begin{equation}
 \operatorname{rank}\!\begin{pmatrix}L_6\\ \kappa_U\end{pmatrix}=7
 \quad\Longleftrightarrow\quad \sigma_U>0.
 \label{eq:p3uscalerank}
\end{equation}
Thus even a nonzero unit differential is insufficient when it lies in the
source-metric row span of the six previous roles.  Physical activation needs
protected-square occupation, quotient descent of $\kappa_U$, and the positive
certificate~\eqref{eq:p3uscaleSchur}.

The unchanged source inventory cannot supply that activation.  The derived
source arrows form a finite $S_4$ relabeling groupoid with zero continuous
tangent dimension, while the occupied CHDF seed orbit is locally isolated
with zero on-shell seed-moduli tangent.  Therefore
\begin{equation}
 T_{x_*}X_{\rm phys}^{\rm curr}=0
 \quad\Longrightarrow\quad
 d\log U_R|_{T_{x_*}X_{\rm phys}^{\rm curr}}=0.
 \label{eq:p3currentsourcenogo}
\end{equation}
The available off-shell seed--body response germ is not a selected family of
physical source arrows and cannot replace this missing tangent.

The adopted MCSCL forward completion provides a conditional reopening.  Its
common calibration tangent $c_U$ satisfies
\begin{equation}
 D\log\ell_H(c_U)=1,\qquad
 D\log\mathcal A_*(c_U)=2,
 \qquad\text{hence}\qquad \kappa_U(c_U)=1
 \label{eq:p3mcsclactivation}
\end{equation}
on the common CHDF/PRUB chart.  Calibration-basicness of the six raw rows is
not an additional physical premise.  For an arbitrary descended six-row
family define its measured common-scale loading and relative part by
\begin{equation}
 v_U:=L_6c_U,\qquad L_6^{\rm rel}:=L_6-v_U\kappa_U.
 \label{eq:p3mcsclseparation}
\end{equation}
Then $L_6^{\rm rel}c_U=0$ and
\begin{equation}
 \begin{pmatrix}L_6^{\rm rel}\\\kappa_U\end{pmatrix}
 =
 \begin{pmatrix}I_6&-v_U\\0&1\end{pmatrix}
 \begin{pmatrix}L_6\\\kappa_U\end{pmatrix}.
 \label{eq:p3scalehortransform}
\end{equation}
The block-triangular row transformation is source-determined, invertible and
has determinant one.  It preserves the complete joint kernel and rank, while
the source-metric Gram changes only by congruence.  Since the relative rows
annihilate $c_U$ but $\kappa_U(c_U)=1$,
\begin{equation}
 \operatorname{rank}\!\begin{pmatrix}L_6\\\kappa_U\end{pmatrix}=7
 \quad\Longleftrightarrow\quad
 \operatorname{rank}L_6^{\rm rel}=6
 \quad\Longleftrightarrow\quad
 L_6^{\rm rel}H_{\rm src}^{-1}(L_6^{\rm rel})^*>0.
 \label{eq:p3scalehorrank}
\end{equation}
Thus the common scale is counted exactly once without discarding it.  This
removes raw calibration-basicness as a completion assumption, but does not
manufacture a missing role direction.  The unrestricted physical parent must
still own the MCSCL character, the protected relation-unit square, the six
occupied evaluation jets, and their positive rank-six relative Gram.

The onto dependency response already supplies the underlying relative
carrier.  Let
$D_{\rm dep}:X_{\rm src}\to Z_{\rm dep}$ be the surjective seven-dimensional
map above and assume quotient descent
\begin{equation}
 \kappa_U=\bar\kappa_UD_{\rm dep},\qquad
 \bar\kappa_U(D_{\rm dep}c_U)=1.
 \label{eq:p3relsplit}
\end{equation}
With $z_U=D_{\rm dep}c_U$, define
\begin{equation}
 E_{\rm dep}^{\rm rel}:=\ker\bar\kappa_U,\qquad
 P_{\rm rel}:=I-z_U\bar\kappa_U.
 \label{eq:p3relprojector}
\end{equation}
Direct multiplication gives
\begin{equation}
 P_{\rm rel}^2=P_{\rm rel},\quad
 \operatorname{im}P_{\rm rel}=E_{\rm dep}^{\rm rel},\quad
 \ker P_{\rm rel}=\operatorname{span}\{z_U\},\quad
 \dim E_{\rm dep}^{\rm rel}=6.
 \label{eq:p3relcarrier}
\end{equation}
Restriction of the positive dependency metric to this kernel is positive
definite.  Hence no additional six-dimensional carrier or metric is needed.
What remains unselected is the occupied typed isomorphism
$T_{\rm rel}:E_{\rm dep}^{\rm rel}\to
E_{P3}^{\ne U}=\operatorname{span}\{B,O,R,Q,M,G\}$.
Orthogonal frames of the same relative carrier have identical upstream
status but different semantic assignments, so positivity alone cannot choose
$T_{\rm rel}$.

The existing typed support ledger reduces, but does not remove, this
ambiguity.  On the non-unit role set retain the named supports
\begin{equation}
 H_{P4}=\{B,M,G\},\qquad
 H_{OI}=\{O,Q\},\qquad
 H_{PM}=\{R,M,G\}.
 \label{eq:p3incidencesupports}
\end{equation}
Here $H_{P4}$ is the non-unit, non-$I$ restriction of the maximal P4
support, while $H_{OI}$ and $H_{PM}$ are the observer--instrument and
photon--matter relation supports.  The incidence fingerprints with respect
to this ordered support family are
\begin{equation}
\begin{array}{c|cccccc}
a&B&O&R&Q&M&G\\ \hline
\chi(a)&100&010&001&010&101&101.
\end{array}
\label{eq:p3incidencefingerprints}
\end{equation}
Therefore the permutation stabilizer of every named support is exactly
\begin{equation}
 \operatorname{Aut}(\mathcal H_{P3}^{\rm rel})
 \simeq S_{\{O,Q\}}\times S_{\{M,G\}},
 \qquad |\operatorname{Aut}(\mathcal H_{P3}^{\rm rel})|=4.
 \label{eq:p3incidenceautomorphism}
\end{equation}
Indeed, $B$ and $R$ have unique fingerprints; the only repeated columns are
the $O/Q$ and $M/G$ pairs, and either pair swap preserves all three supports.
Consequently, if these supports descend to six occupied atomic rays of
$E_{\rm dep}^{\rm rel}$, support preservation reduces the semantic label
freedom to two binary choices.  A source-owned distinction on each unresolved
pair is then necessary and sufficient to make the permutation stabilizer
trivial; one common source object may supply both distinctions.  This is not
yet a construction of $T_{\rm rel}$: the current PS1
action proves neither the atomic support descent nor those two oriented
separators.  It is a finite sharpening of the missing datum, not permission
to select the pair orientations from terminal behavior.

This remainder has an exact quotient formulation.  Put

\begin{equation}
 q_{\rm rel}:=P_{\rm rel}D_{\rm dep}:X_{\rm src}
 \longrightarrow E_{\rm dep}^{\rm rel}.
 \label{eq:p3relativequotient}
\end{equation}
Let $\Xi_{\rm src}:X_{\rm src}\to V_{\rm sep}$ be one source-owned
descriptor.  Since $q_{\rm rel}$ is onto, there is a unique descended map
$\bar\Xi:E_{\rm dep}^{\rm rel}\to V_{\rm sep}$ satisfying
$\Xi_{\rm src}=\bar\Xi q_{\rm rel}$ exactly when
\begin{equation}
 \ker q_{\rm rel}\subseteq\ker\Xi_{\rm src}.
 \label{eq:p3selectordescent}
\end{equation}

\begin{proposition}[Relative pair-selector criterion]
\label{prop:p3pairselector}
Assume that the supports in Eq.~\eqref{eq:p3incidencesupports} descend to
occupied atomic rays with source-owned normalized representatives
$r_a\in E_{\rm dep}^{\rm rel}$.  The subgroup that
preserves both those supports and the values
$\xi_a:=\bar\Xi(r_a)$ is trivial if and only if
\begin{equation}
 \boxed{\xi_O\ne\xi_Q\quad\hbox{and}\quad\xi_M\ne\xi_G.}
 \label{eq:p3pairseparation}
\end{equation}
Thus a single common source descriptor can remove both residual binary
ambiguities without a new carrier or role-dependent gain.
\end{proposition}

\begin{proof}
Equation~\eqref{eq:p3incidenceautomorphism} shows that the support stabilizer
is generated by the independent swaps $(O\ Q)$ and $(M\ G)$.  The first
preserves the descriptor exactly when $\xi_O=\xi_Q$, and the second exactly
when $\xi_M=\xi_G$.  Their common stabilizer is therefore trivial precisely
under Eq.~\eqref{eq:p3pairseparation}.  The descent statement is the universal
property of the onto quotient $q_{\rm rel}$.
\end{proof}

The irreducible observer--source incidence and post-body variational ownership
have the right shape for a two-component candidate: one could distinguish
$O/Q$, the other $M/G$.  Existing rank and terminal-variation theorems do not,
however, prove that they form one source-owned descriptor, satisfy
Eq.~\eqref{eq:p3selectordescent}, or obey
Eq.~\eqref{eq:p3pairseparation}.  Terminal meanings cannot supply those
parent-side facts circularly.  The finite certificate checks the order-four,
one-pair order-two, both-pair trivial and nonfactorizing-null controls; it does
not assert physical base--seed occupation of $\Xi_{\rm src}$.

There is a further exact restriction on the natural candidate. Let
$\Delta_{\rm rec}$ be the present commutative special dagger--Frobenius copy,
$S$ the exchange of its two tensor legs and $P_-=(I-S)/2$. On its copy basis,
\begin{equation}
 \Delta_{\rm rec}(e_k)=e_k\otimes e_k,
 \qquad S\Delta_{\rm rec}=\Delta_{\rm rec},
 \qquad P_-\Delta_{\rm rec}=0.
 \label{eq:p3frobeniusorientation}
\end{equation}

\begin{proposition}[Current-candidate no-go and ordered-port target]
\label{prop:p3currentcandidatenogo}
Relative to the source objects presently constructed, irreducible
observer--source incidence, post-body matter/geometry variational typing and
the commutative Frobenius copy do not define a noncircular
$\Xi_{\rm src}$ satisfying Proposition~\ref{prop:p3pairselector}. A
sufficient finite replacement is one source-owned ordered typed port
incidence
\begin{equation}
 \Pi_{\rm port}:X_{\rm src}\longrightarrow V_{OQ}\oplus V_{MG},
 \qquad
 \Pi_{\rm port}=\bar\Pi_{\rm port}q_{\rm rel},
 \label{eq:p3portincidence}
\end{equation}
whose descended values distinguish both $O/Q$ and $M/G$.
\end{proposition}

\begin{proof}
The observer--source construction supplies a transverse image and rank, but
turning it into an $O/Q$ value requires an additional source-owned codomain
functional or port assignment. Matter/geometry variational ownership is
typed only after the semantic P3 roles have been identified; pulling it back
to $E_{\rm dep}^{\rm rel}$ through the unknown $T_{\rm rel}$ would use the
desired identification in its own construction. Equation
\eqref{eq:p3frobeniusorientation} shows that the current copy has no
exchange-odd component. The current objects therefore coexist with both
residual pair swaps and do not imply Eq.~\eqref{eq:p3pairseparation}.
Conversely, Eq.~\eqref{eq:p3portincidence} obeys the quotient descent
condition by construction, and Proposition~\ref{prop:p3pairselector} removes
the complete residual stabilizer when its two pairwise inequalities hold.
\end{proof}

The previously derived exchange-odd charged-wall operator
$D_{12}=Y_A\otimes I-I\otimes Y_A$ is a type witness for such orientation,
not this selector. It acts on a conditional charged two-leg realization, the
current parent inventory does not own its required $E_{60}$ carrier, and no
factorization through $q_{\rm rel}$ has been established. Proposition
\ref{prop:p3currentcandidatenogo} is therefore a no-go for the current source
inventory, not a theorem that an ordered port incidence cannot exist. Its
physical base--seed occupation remains open.

The existing universal-seed orientation is only one binary sign. This gives
an additional finite no-go.

\begin{proposition}[Single-seed-sign no-go]
\label{prop:p3singleseedbitnogo}
Let
$G_{\rm pair}=S_{\{O,Q\}}\times S_{\{M,G\}}
\simeq\mathbb Z_2\times\mathbb Z_2$ be the residual support stabilizer.
No equivariant action of $G_{\rm pair}$ on one binary seed-orientation line
is faithful. Thus the existing scalar seed sign alone cannot orient both
$O/Q$ and $M/G$.
\end{proposition}

\begin{proof}
Such an action is a character
$\chi:G_{\rm pair}\to\{+1,-1\}\simeq\mathbb Z_2$. Its image has order at
most two, whereas $|G_{\rm pair}|=4$. The first isomorphism theorem therefore
gives $|\ker\chi|\geq2$. At least one nontrivial pair swap remains invisible.
\end{proof}

This result does not multiply the universal morphology seed. One structured
seed or joint base--seed object may carry the faithful bidegree
$(\chi_{OQ},\chi_{MG})$ with trivial common kernel. What is excluded is only
identifying the single global causal/time orientation bit with the complete
relative-role selector. The two typed characters, or an equivalent ordered
port incidence, still require a source derivation and quotient descent.

The required two-character object nevertheless has a canonical finite
representation once its source ownership is supplied. Let $K_{\rm rel}>0$
be the relative metric and let $J_{OQ},J_{MG}$ be commuting
$K_{\rm rel}$-isometric involutions. Assume each has a source-oriented
one-dimensional negative eigenspace, spanned by unit vectors $c_{OQ}$ and
$c_{MG}$, with each anchor even under the other involution. Define
\begin{equation}
 \bar\Pi_{\rm port}(v)=
 \bigl(K_{\rm rel}(c_{OQ},v),K_{\rm rel}(c_{MG},v)\bigr).
 \label{eq:p3portmetricdual}
\end{equation}

\begin{proposition}[Ordered-port representation]
\label{prop:p3portrepresentation}
Suppose $J_{OQ}$ exchanges $r_O,r_Q$ and fixes $r_M,r_G$, whereas
$J_{MG}$ exchanges $r_M,r_G$ and fixes $r_O,r_Q$. If
\begin{equation}
 K_{\rm rel}(c_{OQ},r_O)\ne0,
 \qquad K_{\rm rel}(c_{MG},r_M)\ne0,
 \label{eq:p3portnonzero}
\end{equation}
then Eq.~\eqref{eq:p3portmetricdual} separates both residual pairs and has
trivial support stabilizer. If the raw involutions and anchors preserve
$\ker q_{\rm rel}$, its pullback through $q_{\rm rel}$ is the source
descriptor required by Proposition~\ref{prop:p3pairselector}.
\end{proposition}

\begin{proof}
Isometry and oddness give
$K_{\rm rel}(c_{OQ},r_Q)=
K_{\rm rel}(c_{OQ},J_{OQ}r_O)=
-K_{\rm rel}(c_{OQ},r_O)$, and the first condition in
Eq.~\eqref{eq:p3portnonzero} makes these values distinct. The same argument
applies to $M/G$. Proposition~\ref{prop:p3pairselector} then gives the
trivial stabilizer, while preservation of $\ker q_{\rm rel}$ is precisely
the quotient-descent condition.
\end{proof}

The conditional oriented observer--source double cover already present in
the hub is a candidate for the first odd line, and the once-counted
matter/geometry variational split motivates the second. The current physical
source ledger proves neither joint occupation of those two oriented lines nor
their preservation of $\ker q_{\rm rel}$. Proposition
\ref{prop:p3portrepresentation} is therefore a finite sufficient realization
theorem, not unrestricted parent-law realization.

The two unresolved pairs nevertheless share one variational structure in the
standard leaf. A normalized observer effect evaluates a quantum state,
whereas a metric variation evaluates matter stress:
\begin{equation}
 \mathfrak e_{OQ}(E,\rho)=\operatorname{Tr}(\rho E),
 \qquad
 \mathfrak e_{MG}(h,T)
 =-\frac12\int d^4x\sqrt{-g}\,T^{\mu\nu}h_{\mu\nu}.
 \label{eq:p3dualevaluations}
\end{equation}
These terminal identities do not derive a parent carrier. They do require
every differentiable completion that recovers them to retain two typed
source/response dualities. Let $i\in\{OQ,MG\}$, let
$W_i\subset E_{\rm dep}^{\rm rel}$ be the protected pair subspace, let $h_i>0$
be its inherited metric, let
$\bar\eta_i\in(E_{\rm dep}^{\rm rel})^*$ be a nonzero covector supported on
$W_i$, and admit one conjugate response $v_i\in W_i$. Assume a source-owned
variance representation maps the source and response legs to the two occupied
role rays of the pair. Consider
\begin{equation}
 \Gamma_{\rm port}^{\rm CTC}
 =\sum_i\left[\Psi_i(v_i)+\Psi_i^*(\bar\eta_i)-\bar\eta_i(v_i)\right]
 =\frac12\sum_i\|v_i-\bar\eta_i^\sharp\|_{h_i}^2,
 \qquad \Psi_i(v)=\frac12\|v\|_{h_i}^2.
 \label{eq:p3bicotangent}
\end{equation}

\begin{proposition}[Minimal bicotangent port completion and selection boundary]
\label{prop:p3bicotangent}
In the local positive quadratic one-copy class with inherited metrics and one
common action unit, Eq.~\eqref{eq:p3bicotangent} is the unique completion that
has the exact conjugate graphs $v_i=\bar\eta_i^\sharp$, zero minimized source
remainder and no pair-dependent gain. Its two source/response exchanges are
commuting isometric involutions with one-dimensional odd contrast lines. If
the source pullbacks $\eta_i^{\rm src}=\bar\eta_iq_{\rm rel}$ are nonzero and
the declared variance representation reaches the two role rays, those lines
give the descriptor of Proposition~\ref{prop:p3portrepresentation}.

The frozen current base--seed/body reduct does not entail this completion or
physical occupation of both odd lines.
\end{proposition}

\begin{proof}
Fenchel--Young gives each term in Eq.~\eqref{eq:p3bicotangent} as a
nonnegative square, with equality exactly at $v_i=\bar\eta_i^\sharp$. In a
metric-normalized one-mode chart its joint Hessian is
\begin{equation}
 H_i=\begin{pmatrix}1&-1\\-1&1\end{pmatrix},
\end{equation}
whose response Schur complement is zero. For a general normalized quadratic
candidate with coefficients $(a_i,b_i,c_i)$, exact conjugacy gives
$c_i=a_i$, zero Schur remainder gives $b_i=c_i^2/a_i=a_i$, and the inherited
vertical curvature gives $a_i=1$. Hence the completion is unique in the
declared class. Leg exchange is an isometry of $H_i$; the two exchanges act
on disjoint blocks and commute, and their odd lines are generated by the two
source--response differences. Source basicness is precisely the quotient
descent condition, so Proposition~\ref{prop:p3portrepresentation} applies.

For non-entailment, the hub's cotangent-role control contains two models with
the same base, universal seed, stabilized body, protected covector, metric
and action unit: one omits the conjugate response and one admits it. Their
response ranks differ. Independently, an endpoint-directed source and an
oriented double-cover presentation can share the same local functional, and
the endpoint $\mathbb Z_2^{OS}$ is not the P3 $O/Q$ swap without a typed
intertwiner. Products of these controls have identical frozen reducts but
different ordered-port content. The current reduct therefore cannot select
the bicotangent packet. This is an unrestricted parent-law non-entailment, not a defect
of the sufficient completion.
\end{proof}

Thus standard Born and Hilbert-stress recovery select
Eq.~\eqref{eq:p3bicotangent} as the weakest source-preserving completion in
the declared local class. They do not prove that the earlier physical
base--seed reduct admits and occupies both quotient-basic covectors or owns
their variance representation on the P3 rays. No further port coefficient or
functional-shape search is licensed; the open question is cotangent-role and
typed-port admission by the enriched parent.

The hub's exterior and record results reduce the arbitrariness of that handoff.
Let
\begin{equation}
\mathcal F_{P3}=\Lambda^\bullet E_{P3}\otimes\mathcal R_2,
\qquad \dim\mathcal R_2=2,
\label{eq:p3fockrecord}
\end{equation}
where $\mathcal R_2$ is the already selected two-leg dagger--Frobenius record
carrier.  On $\Lambda^\bullet E_{P3}$ define
\begin{equation}
\mathsf C_{\rm ext}(v)=\varepsilon(v)+\iota_{K_*(v,\cdot)}.
\label{eq:p3extclifford}
\end{equation}
The canonical exterior identities give
$\mathsf C_{\rm ext}(v)^2=K_*(v,v)I$, and tensoring with
$I_{\mathcal R_2}$ leaves this law unchanged.  The exterior carrier has
dimension $2^7=128$ with equal central sectors of dimension $64$; the record
leg doubles them to $(128,128)$.  Since each simple $M_8(\mathbb C)$ module has
dimension eight, the resulting multiplicities are $(16,16)$ and the normalized
trace weights are $(1/2,1/2)$.

This is a functorial, coefficient-free realization of the premise of
Proposition~\ref{prop:p3clifford}; it does not identify the seven Paper-IV
dependency hyperedges with the seven P3 roles.  Its physical premise is instead
joint source ownership and occupation of the already typed $E_{P3}$ role
quotient, its positive metric and the two-leg record carrier.  Those ingredients
exist inside the declared source-irredundant enriched MVP.  More precisely,
if $r_*=\Delta_\tau e_*$ is its normalized occupied Frobenius record state,
then
\begin{equation}
J_{P3,R}\psi=\psi\otimes r_*,\qquad
J_{P3,R}^\dagger J_{P3,R}=I,
\label{eq:p3recordlift}
\end{equation}
and
$(\mathsf C_{\rm ext}(v)\otimes I)J_{P3,R}
=J_{P3,R}\mathsf C_{\rm ext}(v)$.
Thus the common-source lift cannot annihilate occupied P3 support and adds no
gain: joint assembly is closed inside that declared MVP carrier class.  This
conclusion is not a consequence of marginal occupation alone.  On
$\mathbb C^2$, the orthogonal projectors
$P=\lvert0\rangle\!\langle0\rvert$ and
$R=\lvert1\rangle\!\langle1\rvert$ both act nontrivially on
$(\lvert0\rangle+\lvert1\rangle)/\sqrt2$, while $PR=0$.
Consequently the unrestricted base--seed reduct still does not imply that its
physical parent law realizes and occupies the common enriched MVP carrier, but no additional internal
joint-occupation gate remains after that carrier is fixed.

The same packet removes the variance-representation premise left explicit in
Proposition~\ref{prop:p3bicotangent}.  Define
\begin{equation}
 \Phi_{P3}(v)=\mathsf C_{\rm ext}(v)\otimes I_{\mathcal R_2},
 \qquad
 \omega_{P3,R}(A)={1\over256}\operatorname{Tr}(A).
 \label{eq:p3variancerepresentation}
\end{equation}

\begin{proposition}[Canonical port variance and operational occupation]
\label{prop:p3portoccupation}
On the occupied typed P3 exterior--record packet,
\begin{equation}
 \omega_{P3,R}(\Phi_{P3}(v))=0,
 \qquad
 \operatorname{Re}\omega_{P3,R}
 \left(\Phi_{P3}(u)^*\Phi_{P3}(v)\right)=K_*(u,v).
 \label{eq:p3variancegram}
\end{equation}
Hence $\Phi_{P3}$ is a faithful metric-isometric variance representation and
its restrictions to the $O/Q$ and $M/G$ pair subspaces supply the
representation assumed in Proposition~\ref{prop:p3bicotangent} without a new
gain or carrier.

Assume in addition complete degree-one factorization through the typed P3
quotient.  A normalized POVM outcome family induces a nonzero $O/Q$ first
jet.  On every test region where the renormalized matter stress is a nonzero
distribution, metric variation induces a nonzero $M/G$ first jet.  Their
quotient covectors are basic, and the positive Riesz maps give unique nonzero
bicotangent responses.  Thus no further internal port-occupation bit is
required on such support.
\end{proposition}

\begin{proof}
The operators exchange the even and odd exterior grades, so their normalized
trace is zero.  The exterior Clifford relation gives
\begin{equation}
 \Phi_{P3}(u)\Phi_{P3}(v)+\Phi_{P3}(v)\Phi_{P3}(u)
 =2K_*(u,v)I.
\end{equation}
The operators are self-adjoint; normalized-trace cyclicity proves
Eq.~\eqref{eq:p3variancegram}.  Setting $u=v$ gives
$\|\Phi_{P3}(v)\|_{2,\omega}^2=K_*(v,v)$, so positivity of $K_*$ proves
injectivity and nonzero images for all typed role rays.

For the state-dependent claim, POVM normalization gives
$\sum_a\operatorname{Tr}(\rho_OE_a)=1$, so at least one typed effect-source
derivative is nonzero.  A nonzero stress distribution is exactly a nonzero
linear functional on compactly supported metric test variations, hence some
$h$ has
$-\frac12\int\sqrt{-g}\,T^{\mu\nu}h_{\mu\nu}\ne0$.
Complete typed factorization descends both jets to nonzero quotient
covectors; their source pullbacks annihilate the quotient kernel.  Since a
positive metric has an invertible Riesz map, both conjugate responses are
nonzero. QED.
\end{proof}

The nonzero-stress hypothesis is necessary.  An identically zero
renormalized stress on a region gives a zero $M/G$ first jet even though the
abstract $M$ and $G$ role rays and their faithful representation remain
present.  Proposition~\ref{prop:p3portoccupation} therefore closes an
internal representation/activation issue of the already occupied enriched
MVP; it does not prove that the unrestricted physical base--seed parent
selects the typed quotient, the Clifford--record packet or a nonzero-stress
observer support.

The same construction fixes the representation class rather than merely its
dimension.  Since
$\operatorname{Cl}_7(\mathbb C)\simeq M_8(\mathbb C)\oplus M_8(\mathbb C)$
has dimension 128, its left regular module contains eight copies of each
simple eight-dimensional module.  These are exactly the multiplicities of
$\Lambda^\bullet E_{P3}$.  The exterior vacuum has full Clifford-word orbit
rank 128; after the two central record projectors are included, the joint
orbit rank is 256.  Hence the exterior--record packet is a faithful cyclic
regular representation of its P3--record algebra.  Its represented normalized
trace fixes the equal Clifford-sector weights, so no separate cyclicity, GNS
carrier or central-weight choice remains on this subpacket.

\section{Complete-incidence characterization}
\label{app:p3certificate}

This does not
establish minimal-GNS faithfulness or differential zero cokernel for the full
$S0$--$S7$ source signature.  The positive connected extension
$\Gamma_{\rm ext}(z,w)=\Gamma_{P3,R}(z)+h\lVert w-bz\rVert^2/2$ has a
generically nonzero raw coordinate while exact elimination gives
$\operatorname{Schur}_wH_{\rm ext}=H_{P3,R}$.  It therefore disproves raw
coordinate exhaustivity.  It is not, however, a distinct physical-source
completion under the adopted source ontology: in the defect coordinate
$u=w-bz$ it is the centered positive auxiliary
$\Gamma_{P3,R}(z)+h\lVert u\rVert^2/2$ with $u_*=0$, zero independent load and
Schur-neutral elimination.  Physical exhaustivity modulo presentation
auxiliaries is instead reduced to uniqueness of the occupied stable branch
and zero cokernel of the complete typed source incidence.  Those two premises
remain open for the unrestricted physical base--seed parent.  On the declared
source-frozen finite MVP action, strict positivity already supplies the unique
stationary branch, so its sole remaining internal exhaustivity test is the
typed zero-cokernel condition.  Physical base--seed realization and occupation
of that MVP action remain
an upstream physical assumption.

The preceding construction can nevertheless be compressed into one exact
source certificate.  Let $P_{P3}^1$ be the physical degree-one source
quotient, with its declared seven typed roles, and let
\begin{equation}
 \mathcal C_{P3}^{\rm phys}:P_{P3}^1\longrightarrow
 \mathcal A_{\rm body}^{\rm sa}
 \label{eq:p3physicalcertificate}
\end{equation}
be a source-owned incidence into the self-adjoint part of the occupied body
algebra.  Call Eq.~\eqref{eq:p3physicalcertificate} \emph{complete} when
(i) its quotient rank is seven and its kernel is exactly the complete P3
physical null, with no terminal-visible bypass; (ii) after polar
normalization it obeys
\begin{equation}
 \mathcal C(v)^2=K_*(v,v)I;
 \label{eq:p3physicalsquare}
\end{equation}
and (iii) the canonical exterior extension and the two-leg
dagger--Frobenius record algebra used in Eq.~\eqref{eq:p3recordlift} are
jointly occupied on the same source carrier by a cyclic vector with full
algebra orbit (rank 256 in the present finite realization).
Here \emph{minimal} means generated by the physical degree-one source
quotient with no additional P3 bypass; it does not mean smallest Hilbert-space
dimension.  The declared recovery class uses the canonical cyclic regular
exterior--record realization.

\begin{theorem}[Minimal physical P3 certificate]
\label{thm:p3physicalcertificate}
Within the finite source-irredundant Clifford--recovery class, the physical
parent contains the minimal faithful typed P3 Clifford--record subpacket,
uniquely up to typed unitary equivalence, if and only if it owns a complete
incidence \eqref{eq:p3physicalcertificate}.
\end{theorem}

\begin{proof}
For sufficiency, the exact-null condition descends the incidence faithfully
to the seven-dimensional typed quotient.  Polar normalization fixes the
positive metric and Eq.~\eqref{eq:p3physicalsquare}; polarization then gives
the Clifford anticommutator.  The exterior construction
\eqref{eq:p3extclifford}, the occupied record lift
\eqref{eq:p3recordlift}, and Proposition~\ref{prop:p3portoccupation} provide
the faithful cyclic Clifford--record representation and its normalized
trace.  Source irredundancy removes invisible degree-one copies, while
complete-null/no-bypass removes a second terminal-visible realization.
The remaining freedom is therefore typed unitary equivalence.

Conversely, restrict a minimal faithful typed P3 Clifford--record subpacket
to its grade-one self-adjoint Clifford generators and pull them back along
the physical source first jet.  Faithfulness gives rank seven, the physical
quotient gives the exact kernel and no-bypass property, the Clifford law
gives Eq.~\eqref{eq:p3physicalsquare}, and regular cyclic occupation supplies
the same full exterior--record orbit.  Hence the pulled-back incidence is
complete. QED.
\end{proof}
\section{Enriched variational ownership and ambient-extension control}
\label{app:p3ownership}

The ownership question can be closed inside an explicit ordered enriched
completion without fixing the Clifford map in advance.  After the
morphological body has been solved and frozen, vary
$C\in\operatorname{Hom}_{\mathbb R}(E_{P3}^{\rm phys},
\mathcal A_{P3}^{\rm sa})$ and one operator response $S_{P3}$.  For the
source-frozen occupied descriptor $e_{P3}$ define
\begin{align}
 \Gamma_{P3}^{\rm enr}[C,S_{P3};e_{P3}]
 ={}&{\mathcal A_*\over2}\tau\!\left[
 (S_{P3}-C(e_{P3}))^\dagger(S_{P3}-C(e_{P3}))\right]
 \nonumber\\
 &+{\mathcal A_*\over2}\int_{S(E_{P3},K_*)}
 \tau\!\left[\mathfrak R_C(v)^\dagger\mathfrak R_C(v)\right]d\mu(v),
 \qquad
 \mathfrak R_C(v)=C(v)^2-K_*(v,v)I.
 \label{eq:p3enrichedparentaction}
\end{align}
This is a staged parent action: body variables are not varied in the
post-body P3 stage.  The two rows use the same faithful trace and action unit,
and the second counts the primitive operator relation once.

\begin{theorem}[Enriched-parent variational ownership]
\label{thm:p3enrichedownership}
In the declared source-irredundant enriched completion, supplemented by the
complete physical P3 null, cyclic exterior--record occupation and no-bypass
conditions, Eq.~\eqref{eq:p3enrichedparentaction} owns and selects the minimal
faithful P3 Clifford--record subpacket.  Its zero locus is
\begin{equation}
 S_{P3}=C(e_{P3}),
 \qquad
 C(u)C(v)+C(v)C(u)=2K_*(u,v)I,
 \label{eq:p3enrichedzerolocus}
\end{equation}
and exact elimination of $S_{P3}$ leaves
$\mathcal V_{\rm OS}[C]$.  Hence the handoff and source selector are not
double-counted.
\end{theorem}

\begin{proof}
Both terms in Eq.~\eqref{eq:p3enrichedparentaction} are nonnegative.  The
canonical exterior Clifford incidence together with
$S_{P3}=C(e_{P3})$ reaches zero.  Conversely, every global zero has vanishing
graph defect and vanishing operator-square integral.  The first gives the
response equation in Eq.~\eqref{eq:p3enrichedzerolocus}; faithfulness,
continuity and homogeneity give
$C(v)^2=K_*(v,v)I$ for every $v$, and polarization gives the displayed
anticommutator.  Minimization over $S_{P3}$ removes exactly the graph square
and leaves Eq.~\eqref{eq:p3operatorsquarefunctional}.  Theorem
\ref{thm:p3physicalcertificate} then supplies typed-unitary uniqueness of the
minimal occupied subpacket. QED.
\end{proof}

The construction is not an arbitrary terminal repair.  Its typed domain is
the already fixed P3 quotient; its target is the canonical generated exterior
algebra; its metric, trace and unit are inherited; body-first staging forbids
readout-to-body backreaction; and it replaces rather than duplicates the
fixed graph-normal handoff.  Numerically, the graph term minimizes to zero
for both the Clifford and commuting controls, while the reduced enriched
costs remain $0$ and $4/3$, respectively; the zero incidence has normalized
cost one.

This is a conditional source-law closure, not an unrestricted
parent-realization theorem.
The current reduct has no source variable typed as a complete rank-seven
operator incidence.  It admits the equal-Gram controls in
Eq.~\eqref{eq:p3operatorsquarecontrol}, so it cannot entail the selector.
Equation~\eqref{eq:p3operatorsquarefunctional} is also not claimed unique
among all possible positive parent functionals; it is the coefficient-free
representative of the declared one-copy orthogonally averaged squared-defect
class.  The theorem proves ownership by the explicitly declared enriched
parent, not that the earlier reduct physically realizes and occupies that
enrichment.

Theorem~\ref{thm:p3physicalcertificate} is a characterization theorem, not a
derivation of certificate existence from the current base--seed reduct.  Nor
can finite readouts establish that this selected subpacket exhausts the
ambient parent.  Indeed, for any visible action $\Gamma(x)$ and finite
terminal stack $F(x)$, set
\begin{equation}
 \widetilde\Gamma(x,w)=\Gamma(x)
 +{q\over2}(w-bx)^2-j(w-bx),\qquad
 \widetilde F(x,w)=F(x),
 \quad q>0,\ j\ne0.
 \label{eq:p3loadedhiddenextension}
\end{equation}
At every visible stationary point $x_*$, the extension has
$w_*=bx_*+j/q$, preserves $F(x_*)$, and satisfies
\begin{equation}
 \operatorname{Schur}_w H_{\widetilde\Gamma}=H_\Gamma.
 \label{eq:p3loadedhiddenschur}
\end{equation}
Thus the finite terminal family cannot distinguish the original parent from
one with an additional independently loaded source component.  In the
invariant defect coordinate $u=w-bx$, Eq.~\eqref{eq:p3loadedhiddenextension}
is $\Gamma(x)+qu^2/2-ju$: it is not an unloaded presentation auxiliary, but
it is also not a second interacting realization of the same connected
universe packet.  The selected connected packet is exhaustive only after the
independent parent-side condition
$\operatorname{coker}J_{\rm phys}=0$ (or an equivalent source law) is proved.

\section{Finite mixed-Gram operational witness}
\label{app:p3rwwitness}

The remaining cokernel is not distributed over all named terminals.  Existing
hub theorems make the P1 gravity--time--radiation/matter differential core and
the finite quantum/P3 restriction separately onto on their declared working
packets.  The unresolved datum is their typed gluing: either
$\ker J_Q\subseteq\ker J_{P1}$ and a source-owned intertwiner
$J_{P1}=\iota_{P1}J_Q$ places P1 inside the complete quantum descriptor, or an
independent P1 factor requires a nonseparable joint incidence.  The relevant
finite calculation is therefore the ROOT-to-QFT preparation/stress
cross-Jacobian, not an additional terminal rank counted by name.
Equation~\eqref{eq:rootp1factor} closes its genuinely quantum stress component
on the declared same-state two-mode sector, and
Eq.~\eqref{eq:rootp1tangentfactor} closes the full P1 restriction on the pure
post-body ROOT tangent.  Equations~\eqref{eq:stagedglue}--
\eqref{eq:stagedgluerank} eliminate that tangent exactly and reduce the
residual kernel test to one operator on the complementary source domain.
The controls in Eq.~\eqref{eq:restcontrols} prove that present marginal data
do not select its rank. They do not decide whether those remaining directions are basic on the
complete quantum descriptor or belong to an additional jointly occupied P1
factor.

The preceding ranks belong to the source-minimal primitive carrier.  They
must not be identified with the row rank of an observer instrument evaluated
on a larger morphology tangent.  The latter can contain several endpoint
functionals of the same primitive role without creating additional parent
source directions.  The following finite witness makes that distinction
explicit and also proves that the mixed-Gram increment can be nonzero.

\subsubsection{Finite operational witness}

\begin{proposition}[Finite protected Robertson--Walker witness]
\label{prop:rwwitness}
On the endpoint-normalized eight-mode homogeneous tangent
\begin{equation}
e_n(t)=t^n-1,\qquad n=1,\ldots,8,
\label{eq:rwbasis}
\end{equation}
take three comoving source times
$t_s\in\{0.21,0.47,0.79\}$ and the metric P1 rows
\begin{align}
Z_s[f]&=-f(t_s),\nonumber\\
J_s[f]&=(1-t_s)f(t_s)-\int_{t_s}^{1}f(t)\,dt,\nonumber\\
E[f]&=-\ddot f(1).
\label{eq:rwmetricrows}
\end{align}
Let the two real ROOT rows be the real and imaginary parts of
\begin{equation}
q_\omega[f]=\int_0^1 e^{-2i\omega t}\dot f(t)\,dt,
\qquad \omega=2.35.
\label{eq:rwquantumrow}
\end{equation}
On the protected quadratic body branch the direct Green-clock third-jet row
vanishes.  The resulting finite matrices satisfy
\begin{equation}
\operatorname{rank}A_{P1}^{RW}=7,
\qquad
\operatorname{rank}\begin{pmatrix}A_{P1}^{RW}\\
\Re q_\omega\\ \Im q_\omega\end{pmatrix}=8,
\qquad r_{Q\perp}^{RW}=1.
\label{eq:rwranks}
\end{equation}
Hence a common-source quantum row can detect one morphology direction erased
by this finite geometric observer packet.
\end{proposition}

\begin{proof}
Substitution of Eq.~\eqref{eq:rwbasis} into
Eq.~\eqref{eq:rwmetricrows} gives three redshift rows, three isotropic Jacobi
rows and one tidal row.  Direct singular-value evaluation gives rank seven
and a one-dimensional null.  Numerical quadrature of
Eq.~\eqref{eq:rwquantumrow} gives a nonzero norm
$5.673092\times10^{-4}$ after projection to this null, so stacking its real
and imaginary parts raises the rank to eight.  The calculation is regenerated
by \path{scripts/audit_theorem_certificates.py}.
\end{proof}

This is a source-matched positive operational witness, not unrestricted
parent-law realization.  It is resolution dependent: sufficiently complete homogeneous
geometric endpoint data can already become injective, in which case the same
quantum functional adds zero rank.  The one-dimensional unresolved fiber is
also not by itself a Hilbert space or a proof of Born structure.  Finally,
Eq.~\eqref{eq:rwranks} does not contradict Eq.~\eqref{eq:deduprank}: the
former is an observer-operator rank on an eight-mode morphology tangent,
whereas the latter is the de-duplicated primitive source-carrier rank.

\section{Class-minimal support--port enriched-law certificate}
\label{app:p3minimalparentlaw}

The support--port reconstruction separates the data that define the proposed
physical law from the conclusions generated by that law.  The support
projectors and common-unit character specify the admissible source domain;
the cotangent, operator-square and record terms specify one staged
once-counted functional; and the decoder specifies noncollapsing body
realization.  They are not independent terminal fits.

\begin{proposition}[Variational occupation, class-relative componentwise minimality and narrow-reduct no-go]
\label{prop:p3minimalparentcertificate}
In the finite source-irredundant one-copy class used in
Theorem~\ref{thm:p3supportportparent}, the normalized P3 functional is
nonnegative and has an explicit simultaneous zero.  Consequently its global
minimum is zero and every global minimizer
lies on the complete packet zero locus; no separate P3 common-zero occupation
assumption is required.  Moreover, every clause of the proposed P3 law is
required for that theorem's conclusion in the following precise sense:
\begin{enumerate}[label=(\roman*),nosep]
\item deleting the support flag leaves the B/R singleton blocks and O/Q,
  M/G planes unidentified;
\item deleting either ordered cotangent port leaves the corresponding
  nontrivial pair swap in the typed stabilizer;
\item placing $z_U$ in the relative span lowers the source rank to six even
  at fixed total row cost;
\item deleting the left-decoder premise permits a rank-six body realization;
\item deleting the operator-square term admits a commuting packet with the
  same normalized marginal trace Gram;
\item deleting the pointed root admits a faithful full-rank nontracial orbit;
  and
\item deleting cyclicity or no bypass permits unoccupied or additional record
  support.
\end{enumerate}
The two cotangent response legs and the graph-normal P3 response can be
eliminated exactly with zero additional source stiffness, so their inclusion
does not duplicate the retained operator-square law.  The unchanged narrow
base--seed reduct admits the first, third and fifth controls with the same
declared marginal upstream data and therefore does not entail
$\mathcal L_{P3}^{\min}$.
\end{proposition}

\begin{proof}
The shifted stabilized-body term is nonnegative on its declared basin.  Each
cotangent contribution is a Fenchel--Young square; the P3 graph and
operator-square contributions are squared faithful-trace norms; and the
record contribution is the normalized squared record defect.  Choose the
canonical exterior Clifford incidence, set its graph response equal to its
source evaluation, use the two metric-dual ordered ports, the stabilized body
and its exact left decoder, and take the occupied empty exterior word with
the normalized two-leg record root.  Every displayed defect then vanishes
simultaneously.  Hence the action floor is zero.  A sum of nonnegative terms
can attain that floor only when every term vanishes, proving variational
occupation of the packet.

The support stabilizer is
$S_{\{O,Q\}}\times S_{\{M,G\}}$.  One port removes one factor and two ports
remove both.  The independent/dependent unit controls have ranks seven and
six with equal total squared row norm.  The redundant body realization has a
left inverse, whereas identifying its unit column with another column lowers
the rank and makes an exact left inverse impossible.  Clifford generators and
the commuting control have the same identity trace Gram, but only the former
have a state-neutral square and vanishing mixed anticommutators.  The
positive noncentral root-density control has full orbit rank and equal central
weights but a nonidentity Gram.  The cyclic/no-bypass controls are the
complete-incidence and pointed-orbit countermodels proved above.  Finally,
each normalized Fenchel--Young port has zero source Schur remainder, and
minimizing the graph square in Eq.~\eqref{eq:p3enrichedparentaction} leaves
the operator-square functional exactly.  These facts prove all clauses and
the same-reduct non-entailment.  QED.
\end{proof}

The proposition is a class-relative minimality certificate for the stated
theorem, not a claim that no richer physical parent law can realize the same
packet.  The superscript ``min'' denotes only componentwise minimality in the
declared finite, source-irredundant one-copy class.  Under the four-dimensional
recovery-completion policy and within that class, the O/Q and M/G dualities,
their de-duplicated relation supports, the common unit and body noncollapse
are the weakest preimages of the already declared standard recovery targets.
The state-neutral operator-square norm is the genuinely new Tau-specific
cross-role clause.  The pointed cyclic root is instead the minimal
record-side response, within the class, to the faithful full-rank nontracial
countermodel.  Hence the paper no longer leaves an unspecified
``complete rank-seven incidence'' as a free input: it proposes the explicit
law in Eq.~\eqref{eq:p3minimalparentlaw}, derives the packet from it, and
derives occupation from global minimization rather than assuming its common
zero separately.  Its
empirical truth and unrestricted ambient-source exhaustivity remain separate
questions.
